\subsection{关于有限个测度乘积的积分}

前面的结果可以毫无困难地推广到有限个测度的乘积。例如，设 \(T_{1}, T_{2}, T_{3}\) 是三个局部紧空间，\(\mu_{i}\) 是 \(T_{i}\) 上的正测度（i = 1, 2, 3），并设 \(\nu = \mu_{1} \otimes \mu_{2} \otimes \mu_{3}\) 是 \(T = T_{1} \times T_{2} \times T_{3}\) 上的乘积测度。设 f 是 \(\nu\) -可积、取值于 \(\overline{R}\) 或一个 Banach 空间中的函数；Lebesgue--Fubini 定理的第一次应用表明，除去点 \((t_{1}, t_{2}) \in T_{1} \times T_{2}\) （组成一个对 \(\mu_{1} \otimes \mu_{2}\) 可忽略的集）外，函数 \(t_{3} \mapsto \mathbf{f}(t_{1}, t_{2}, t_{3})\) 是 \(\mu_{3}\) -可积的，函数

\[
(t _ {1}, t _ {2}) \mapsto \int \mathbf {f} (t _ {1}, t _ {2}, t _ {3}) d \mu_ {3} (t _ {3}),
\]

（在 \(\mathrm{T}_1\times \mathrm{T}_2\) 中几乎处处有定义）是 \((\mu_{1}\otimes \mu_{2})\) -可积的，且

\[
\iiint \mathbf {f} (t _ {1}, t _ {2}, t _ {3}) d \nu (t _ {1}, t _ {2}, t _ {3}) = \iint d \mu_ {1} (t _ {1}) d \mu_ {2} (t _ {2}) \int \mathbf {f} (t _ {1}, t _ {2}, t _ {3}) d \mu_ {3} (t _ {3}).
\]

同一定理的第二次应用表明，对几乎每个 \(t_{1} \in T_{1}\) ，函数 \(t_{2} \mapsto \int \mathbf{f}(t_{1}, t_{2}, t_{3}) d\mu_{3}(t_{3})\) 在 \(T_{2}\) 中几乎处处有定义且是 \(\mu_{2}\) -可积的；此外，函数

\[
t _ {1} \mapsto \int d \mu_ {2} (t _ {2}) \int \mathbf {f} (t _ {1}, t _ {2}, t _ {3}) d \mu_ {3} (t _ {3}),
\]

（在 \(T_{1}\) 中几乎处处有定义）是 \(\mu_{1}\) -可积的，且

\[
\iiint \mathbf {f} (t _ {1}, t _ {2}, t _ {3}) d \nu (t _ {1}, t _ {2}, t _ {3}) = \int d \mu_ {1} (t _ {1}) \int d \mu_ {2} (t _ {2}) \int \mathbf {f} (t _ {1}, t _ {2}, t _ {3}) d \mu_ {3} (t _ {3}).
\]

类似地可以证明，对几乎每个 \(t_{1} \in T_{1}\) ，函数 \((t_{2}, t_{3}) \mapsto \mathbf{f}(t_{1}, t_{2}, t_{3})\) 是 \((\mu_{2} \otimes \mu_{3})\) -可积的，函数

\[
t _ {1} \mapsto \iint \mathbf {f} (t _ {1}, t _ {2}, t _ {3}) d \mu_ {2} (t _ {2}) d \mu_ {3} (t _ {3}),
\]

（几乎处处有定义）是 \(\mu_{1}\) -可积的，且

\[
\iiint \mathbf {f} (t _ {1}, t _ {2}, t _ {3}) d \nu (t _ {1}, t _ {2}, t _ {3}) = \int d \mu_ {1} (t _ {1}) \iint \mathbf {f} (t _ {1}, t _ {2}, t _ {3}) d \mu_ {2} (t _ {2}) d \mu_ {3} (t _ {3}).
\]

我们把对上面就两个测度乘积证明的其他结果作同样推广的任务留给读者。

